\documentclass[11pt]{article}
\usepackage[margin=1in]{geometry}
\usepackage{amsmath,amssymb}
\usepackage{graphicx}
\usepackage{booktabs}
\usepackage{hyperref}
\usepackage{xcolor}
\usepackage{caption}
\usepackage{placeins}
\usepackage[most]{tcolorbox}
\usepackage{algorithm}
\usepackage{algpseudocode}

\newtcolorbox{note}{
  colback=gray!8, colframe=gray!40, boxrule=0.4pt, arc=1pt,
  left=6pt, right=6pt, top=4pt, bottom=4pt,
  fonttitle=\bfseries, title=Note
}

\title{Cross-Sectional Reproduction of the SPX Implied-Vol Smile, the\\
VIX Term Structure, and the VIX-Option Smile by the SP500\\
Feedback-Layer ESN (2026-02-19 Snapshot)}
\author{Companion note to the ESN A2-981003 joint-calibration family\\
(feedback-layer/$\alpha$ version)}
\date{September 7, 2026}

\begin{document}
\maketitle

\begin{abstract}
The joint-calibration note family fits the ESN-A2 architecture to 5
time-series stylised facts of daily returns -- roughness, leverage,
the Taylor effect, the Zumbach effect, and excess kurtosis -- none of
which touch option prices at all. This note asks a different,
cross-sectional question: turned into an option-pricing engine (the
simulator is already a discrete-time exponential martingale under its
own measure -- Proposition~1 in Section~\ref{sec:pricing} -- so no
risk-premium adjustment is needed), does the SP500 calibration
($H=0.386$, \texttt{rough\_scale}$=0.672$, $\lambda_{lo}=10^{-5}$,
$\lambda_{hi}=4.8809$, $m_1=+0.1745$, $\alpha=0.220$) reproduce the
actual SPX implied-vol smile, the VIX term structure, and the
VIX-option smile observed on 2026-02-19 -- and can a genuine
cross-sectional calibration close the gap? As calibrated for the
time series alone, no: the model overstates level by $\approx1.8\times$
on both surfaces (30-day VIX $=36.0$ vs.\ market $=19.6$) and shows a
real shape mismatch on both smiles. Section~\ref{sec:jointcalib} fits
all 8 free parameters ($H,\texttt{rough\_scale},\lambda_{lo},\lambda_{hi},
m_1,m_2,\alpha,\Delta b_0$) jointly against a vega-weighted implied-vol
objective, sharpened toward the money with an explicit boost on
$m\ge1$ points, against \emph{all 6 quoted expiries at once}: SPX near
and next term, plus VIX near, next, 118d, and 153d -- calibrating
directly on the mid of each strike's put and call implied vol (1
point per strike, not 2 independent targets) rather than the raw put
and call quotes separately. The result: SPX near/next weighted RMSE
$0.029$/$0.026$, VIX near/next/118d/153d $0.29$/$0.08$/$0.09$/$0.11$
(confirm-quality, averaged over 2 independent Monte-Carlo seeds --
Table~\ref{tab:calibresult}). The SPX smile now coincides with the
market on both the $m<1$ and $m\ge1$ sides at comparable RMSE
(previously $m\ge1$ was $\approx3.5\times$ worse, now $\approx1.4\times$);
the VIX smile tracks the market closely across most of its quoted
range at 55d and 118d, with a real but smaller residual gap at 27d
and 153d -- an honest limit of the shared 8-parameter architecture
fitting 4 VIX maturities' own term structure at once, not a
fitting failure. Priced at 2 further, entirely held-out SPX
maturities (Section~\ref{sec:heldout}), the fit does not generalize
uniformly: the SPX put wing undershoots increasingly with maturity --
an honest limit of a 2-maturity SPX calibration set, not a
contradiction of the results above. The VIX-option test itself needs
a genuinely different, nested Monte-Carlo algorithm -- detailed in
full in Section~\ref{sec:vix-algo} -- since VIX options price the
model-free variance rate observed at a \emph{future} date, not the
underlying level. All tests use the real
\texttt{esn\_base.py}/\texttt{joint\_calibration.py} engine end to
end, with no modification.
\end{abstract}

\tableofcontents

\section{Market data}
\label{sec:data}
2 option-chain snapshots, both dated 2026-02-19: SPX index options
(11 expiries, 20 days to 5.8 years) and VIX options (4 expiries, 27 to
153 days). Quotes are mid implied volatilities (Black-76 convention),
not dollar premia. The CBOE VIX methodology is applied literally,
\begin{equation}
\sigma^2 = \frac{2}{T}\sum_i \frac{\Delta K_i}{K_i^2}\, e^{RT}Q(K_i)
- \frac{1}{T}\Big[\frac{F}{K_0}-1\Big]^2,
\label{eq:cboe}
\end{equation}
with $e^{RT}Q(K_i)$ built directly from each strike's own mid-IV via
undiscounted Black-76 (call for $K>K_0$, put for $K<K_0$, both
averaged at $K_0$) -- since $Q(K)=e^{-RT}\times\text{Black76}(K)$
exactly under this convention, $R$ never needs to be estimated
separately. Applied to the near-term (2026-03-20, $T_1=29$d) and
next-term (2026-04-17, $T_2=57$d) SPX expiries and interpolated to a
constant 30-day maturity (2-step CBOE procedure), this reproduces the
file's own quoted VIX spot (19.34) to within 1.5\% (19.63 computed),
validating the replication. The identical procedure on the VIX option
chain gives a VVIX-style vol-of-VIX reading of 108.3\% (used only as
context; not part of either test below).

\FloatBarrier
\section{The model}
\label{sec:model}
This section documents the full architecture as implemented in
\texttt{esn\_base.py} -- reservoir, $z$-bank, quadratic channel,
feedback layer, and bias calibration -- at the level needed to
understand exactly what is being priced in Section~\ref{sec:pricing}.

\subsection{Overview}
The simulator is a 2-layer echo-state network driving a single
variance signal, which in turn drives daily log-returns:
\begin{center}
\begin{tabular}{ccccc}
reservoir $r_t$ ($N_r$ OU modes) & $\to$ & pre-activation $\eta_t$ &
$\to$ & $\sigma_t^{\text{raw}}, \sigma_t^{\text{used}}, x_t$ \\
$z$-bank $z_t$ ($N_z$ nonlinear units, reading 2 taps off $r_t$) &
$\nearrow$ & & &
\end{tabular}
\end{center}
Both banks are updated once per trading day ($dt=1$) from a single
shared innovation $\varepsilon_t\sim\mathcal N(0,1)$ -- the same
innovation that both advances the state \emph{and} generates that
day's return, which is the mechanical origin of the leverage effect.

\subsection{The reservoir}
\label{sec:reservoir}
$r_t=(r_{t,1},\dots,r_{t,N_r})$ is a bank of independent
Ornstein--Uhlenbeck modes with geometrically spaced rates,
\begin{equation}
\lambda_i = \lambda_{lo}\Big(\frac{\lambda_{hi}}{\lambda_{lo}}\Big)^{\frac{i-1}{N_r-1}},
\qquad i=1,\dots,N_r,
\label{eq:lambdas}
\end{equation}
updated as
\begin{equation}
r_{t+1,i} = \alpha_i^{\text{OU}} r_{t,i} + c_i\,\varepsilon_t, \qquad
\alpha_i^{\text{OU}} = e^{-\lambda_i dt}, \qquad c_i = \sqrt{1-(\alpha_i^{\text{OU}})^2}.
\label{eq:reservoirupdate}
\end{equation}
The reservoir is read through a single direction $q\in\mathbb R^{N_r}$,
\begin{equation}
q_i \propto \lambda_i^{0.5-H}, \qquad
q \leftarrow \frac{q}{\lVert q\rVert}, \qquad
q \leftarrow \frac{q}{\sqrt{q^\top C q}}, \qquad
C_{ij} = \frac{b_i b_j}{\lambda_i+\lambda_j}, \quad b_i=\sqrt{2\lambda_i}.
\label{eq:qdef}
\end{equation}

\begin{note}
$C$ in Eq.~\eqref{eq:qdef} is the \emph{continuous-time} stationary
covariance of the OU bank ($b_ib_j/(\lambda_i+\lambda_j)$), used only
to normalise $q$ so that $q^\top C q=1$. The simulator itself runs in
\emph{discrete} time with $dt=1$, whose own exact stationary
covariance is $\mathrm{Cov}^{\text{disc}}_{ij}=c_ic_j/(1-\alpha_i^{\text{OU}}\alpha_j^{\text{OU}})$
-- close to $C$ only when $\lambda_i dt\ll1$, which is not guaranteed
here since $\lambda_{hi}$ can reach $5.0$. This is exactly what
\texttt{esn\_base.py} does (not a simplification introduced by this
note); it means $q^\top\mathrm{Cov}^{\text{disc}}q\ne1$ exactly, a
second-order normalisation slack inherited as-is.
\end{note}

\subsection{The $z$-bank}
\label{sec:zbank}
$z_t=(z_{t,1},\dots,z_{t,N_z})$ has its own geometrically spaced decay
rates $a_j=a_{lo}(a_{hi}/a_{lo})^{(j-1)/(N_z-1)}$. Each unit reads 2
fixed reservoir taps, $p_1=r_{t+1,\mathrm{fast}(j)}$,
$p_2=r_{t+1,\mathrm{slow}(j)}$ (fast/slow tap indices defined so that
$\mathrm{fast}(j)$ always sits in the reservoir's upper half and
$\mathrm{slow}(j)$ in its lower half, monotone in $j$):
\begin{equation}
z_{t+1,j} = e^{-a_jdt}z_{t,j} + (1-e^{-a_jdt})\, m_{t+1}\,z_{\text{strength}}\tanh(u_{t,j}),
\label{eq:zupdate}
\end{equation}
\begin{equation}
u_{t,j} = s_jp_1p_2 + e_{\text{strength}}(0.7p_1^2+0.3p_2^2-1) - l_{\text{strength}}(0.7p_1+0.3p_2)
+ zz_jz_{t,j} + \tfrac12 lz_{\text{strength}}(z_{t,j-1}+z_{t,j+1}),
\label{eq:udef}
\end{equation}
with the gate $m_{t+1}=\max(0,\,1-\gamma_{\text{norm}}\lVert r_{t+1}\rVert/\sqrt{N_r})$
suppressing $z$-bank updates whenever the reservoir's current norm is
large, and $s_j\in\{-1,+1\}$, $zz_j\in[-zz_{\text{scale}},zz_{\text{scale}}]$
drawn once from \texttt{MATRIX\_SEED} at construction time
(deterministic given the seed, never re-drawn during simulation). The
$z$-bank enters the pre-activation through a fixed per-mode readout
profile,
\begin{equation}
w_{j,z} = \frac{z_{r,j}}{\sqrt{N_z}}, \qquad
z_{r,j} = zr_{lo}\Big(\frac{zr_{hi}}{zr_{lo}}\Big)^{\frac{j-1}{N_z-1}}.
\label{eq:zreadout}
\end{equation}

\subsection{Quadratic channel}
\label{sec:quad}
A second, independent nonlinear channel reads the reservoir's own
energy and its projection onto $q$:
\begin{equation}
r_t^\top Q r_t = \frac{1}{N_r}\Big(m_1\lVert r_t\rVert^2 + m_2(q^\top r_t)^2\Big),
\qquad Q = \frac{1}{N_r}\big(m_1 I + m_2 qq^\top\big),
\label{eq:quad}
\end{equation}
with $m_2\equiv0$ throughout this note family (only $m_1$ is
calibrated); no positive-definiteness constraint is imposed on $Q$.

\subsection{Pre-activation, bias calibration, and the feedback layer}
\label{sec:feedback}
\begin{equation}
\eta_t = b_0 + \rho\cdot\texttt{rough\_scale}\cdot(q^\top r_t) + w_z^\top z_t
+ \frac{1}{N_r}\big(m_1\lVert r_t\rVert^2+m_2(q^\top r_t)^2\big),
\label{eq:eta}
\end{equation}
where $\rho=\texttt{rough\_orientation}\in\{-1,+1\}$ is an
architecture-level sign flag (fixed at $\rho=-1$ throughout this note
family -- Table~\ref{tab:arch}). The raw pre-activation variance and
its EWMA-smoothed counterpart are
\begin{equation}
(\sigma_t^{\text{raw}})^2 = \text{SIG\_MIN}^2+\text{softplus}(\eta_t)^2,
\qquad
(\sigma_t^{\text{used}})^2 = (1-\alpha)(\sigma_{t-1}^{\text{used}})^2+\alpha(\sigma_t^{\text{raw}})^2,
\label{eq:sigdef}
\end{equation}
with returns generated from the smoothed signal,
\begin{equation}
x_{t+1} = -\tfrac12(\sigma_t^{\text{used}})^2 dt + \sigma_t^{\text{used}}\sqrt{dt}\,\varepsilon_t.
\label{eq:xdef}
\end{equation}
$\alpha=1$ recovers $\sigma_t^{\text{used}}\equiv\sigma_t^{\text{raw}}$
exactly (single-signal model). The bias $b_0$ is \emph{not} a free
parameter -- it is fixed, for every asset, so that the reservoir's own
rest state ($r=0,z=0$) reads exactly a $\text{TV\_ANN}=20\%$
annualised baseline:
\begin{equation}
b_0 = \text{softplus}^{-1}\Big(\sqrt{\max(\text{TV\_DAY}^2-\text{SIG\_MIN}^2,\,10^{-15})}\Big),
\qquad \text{TV\_DAY}=\frac{\text{TV\_ANN}}{\sqrt{252}}.
\label{eq:b0}
\end{equation}

\begin{note}
$m_1=+0.1745>0$ in the SP500 calibration means the quadratic term in
Eq.~\eqref{eq:quad} is strictly non-negative and \emph{always} pushes
$\eta_t$ above the level Eq.~\eqref{eq:b0} was calibrated for; combined
with the convexity of $\mathrm{softplus}(\cdot)^2$, $\mathbb
E[(\sigma_t^{\text{raw}})^2]$ is structurally biased above the 20\%
baseline. Nothing in the 6-parameter joint-calibration objective
(Hurst/leverage/Taylor/Zumbach/kurtosis only) constrains this -- the
resulting average level is a free side-effect, not a fit target. This
is the main driver of the level gap in Section~\ref{sec:results}.
\end{note}

\subsection{Stability condition}
\label{sec:stability}
\texttt{build\_esn\_matrices} also returns
$\kappa_0=\rho\cdot\texttt{rough\_scale}\cdot q^\top\sqrt{2\lambda}$;
the calibration code treats $\kappa_0\ge0$ as an invalid/unstable
parameter region and rejects it outright (a hard constraint enforced
during optimisation, not a soft penalty). For the SP500 row,
$\kappa_0=-0.887<0$: valid.

\subsection{Full architecture and parameter table}
\label{sec:params}
\begin{table}[h]
\centering
\caption{Fixed architecture constants (\texttt{OPT\_ARCH}, identical
across all 9 assets in the joint-calibration note family).}
\label{tab:arch}
\begin{tabular}{lc}
\toprule
$N_r$ & 22 \\
$N_z$ & 29 \\
\texttt{rough\_orientation} $(\rho)$ & $-1.0$ \\
$z_{\text{strength}}$ & 0.27136 \\
$e_{\text{strength}}$ & 3.18426 \\
$l_{\text{strength}}$ & 0.18199 \\
$\gamma_{\text{norm}}$ & 1.30810 \\
$lz_{\text{strength}}$ & 0.06081 \\
$zz_{\text{scale}}$ & 0.03121 \\
$\mathrm{sign\_prob\_neg}$ & 0.22278 \\
\texttt{MATRIX\_SEED} & 202695547565 \\
\midrule
$a_{lo},a_{hi}$ ($z$-bank rates) & $1/400,\ 1/7$ \\
$zr_{lo},zr_{hi}$ ($z$-readout profile) & $0.055,\ 0.25$ \\
$m_2$ & 0 (fixed) \\
$\text{TV\_ANN}$, $\text{SIG\_MIN}$ & $0.20$,\ $0.01/\sqrt{252}$ \\
\bottomrule
\end{tabular}
\end{table}

\begin{table}[h]
\centering
\caption{SP500 calibrated parameters (Table 1, feedback-layer
joint-calibration note) -- the ONLY inputs that vary by asset.}
\label{tab:sp500}
\begin{tabular}{lc}
\toprule
$H$ & 0.386 \\
\texttt{rough\_scale} & 0.672 \\
$\lambda_{lo}$ & $1\times10^{-5}$ \\
$\lambda_{hi}$ & 4.8809 \\
$m_1$ & $+0.1745$ \\
$\alpha$ & 0.220 \\
\bottomrule
\end{tabular}
\end{table}

\subsection{Simulation algorithm}
\label{sec:algo}
\begin{algorithm}[h]
\caption{One simulated day, given state $(r_t,z_t,(\sigma_{t-1}^{\text{used}})^2)$}
\label{alg:day}
\begin{algorithmic}[1]
\State draw $\varepsilon_t\sim\mathcal N(0,1)$
\State $\eta_t \gets$ Eq.~\eqref{eq:eta} using $r_t,z_t$ (state \emph{before} this step's update)
\State $(\sigma_t^{\text{raw}})^2,(\sigma_t^{\text{used}})^2 \gets$ Eq.~\eqref{eq:sigdef}
\State $x_{t+1} \gets$ Eq.~\eqref{eq:xdef}, using $\varepsilon_t$
\State $r_{t+1} \gets \alpha^{\text{OU}}\odot r_t + c\,\varepsilon_t$ \Comment{Eq.~\eqref{eq:reservoirupdate}, SAME $\varepsilon_t$}
\State $m_{t+1} \gets \max(0,1-\gamma_{\text{norm}}\lVert r_{t+1}\rVert/\sqrt{N_r})$
\For{$j=1,\dots,N_z$}
    \State $z_{t+1,j} \gets$ Eq.~\eqref{eq:zupdate}, reading $p_1,p_2$ from $r_{t+1}$
\EndFor
\State \Return $x_{t+1}, r_{t+1}, z_{t+1}, (\sigma_t^{\text{used}})^2$
\end{algorithmic}
\end{algorithm}
Re-using the same $\varepsilon_t$ for both the return (step 4) and the
reservoir update (step 5) is what generates the leverage effect; it is
also, as shown next, exactly what makes the simulator a valid pricing
engine.

\FloatBarrier
\section{Turning the model into a cross-sectional pricing engine}
\label{sec:pricing}

\begin{note}
\textbf{Proposition 1 (martingale property).}
Under Eq.~\eqref{eq:xdef}, conditional on the state at day $t$ (hence
on $\sigma_t^{\text{used}}$, which is $\mathcal F_t$-measurable),
\[
\mathbb E\big[e^{x_{t+1}}\,\big|\,\mathcal F_t\big]
= e^{-\frac12(\sigma_t^{\text{used}})^2}\,\mathbb E\big[e^{\sigma_t^{\text{used}}\varepsilon_t}\big|\mathcal F_t\big]
= e^{-\frac12(\sigma_t^{\text{used}})^2}\,e^{\frac12(\sigma_t^{\text{used}})^2} = 1.
\]
Chaining over $N$ days, $S_t=S_0\exp(\sum x)$ is a discrete-time
exponential martingale under the simulator's own measure. It can
therefore be used directly as a Monte-Carlo pricing engine with
forward $F_{\text{model}}=1$ (moneyness units), with no risk-premium
adjustment or re-weighting.
\end{note}

\subsection{SPX-smile algorithm}
\label{sec:spx-algo}
Turning the simulator into a cross-sectional pricing engine for a
single expiry $T$ (calendar days $N_{\text{cal}}$) proceeds as follows.

\begin{note}
\textbf{ATM-weighted moneyness grid.} Rather than a uniform strike
step, the inversion uses a denser grid near the money and a coarser
one in the wings:
\begin{equation}
\Delta m(m) = \begin{cases}
0.005 & m\in[0.85,1.20] \text{ (SPX)}, \text{ or } 0.01 \text{ (VIX-option)} \\
0.02 & \text{otherwise (SPX)}, \text{ or } 0.04 \text{ (VIX-option)}
\end{cases}
\label{eq:grid}
\end{equation}
This matters for 2 reasons. First, \texttt{cboe\_variance\_from\_grid}
(Eq.~\eqref{eq:cboe}) uses the local strike spacing $\Delta K_i$ as its
own quadrature weight, so a denser ATM grid directly sharpens the
variance integral exactly where the smile's curvature is highest,
without paying the computational cost of that density in the flat
wings. Second, it concentrates the strike-by-strike Black-76 inversion
where real markets are actually liquid and where the comparison
against the true smile is most meaningful -- deep-wing strikes are
priced from very few Monte-Carlo paths reaching that far into the
tail, so their own inverted vols are the noisiest and least
economically important part of the curve.
\end{note}

\begin{algorithm}[H]
\caption{SPX cross-sectional smile at horizon $N_{\text{cal}}$ calendar days}
\label{alg:spx}
\begin{algorithmic}[1]
\State $N \gets \mathrm{round}(N_{\text{cal}}\times252/365)$ \Comment{trading-day equivalent}
\State $T_{\text{years}} \gets N_{\text{cal}}/365$
\State simulate $K$ independent chains, $\mathrm{burn}+N$ days each, from $r_0=z_0=0$, $(\sigma^{\text{used}}_{-1})^2=0$ (Algorithm~\ref{alg:day}, chained)
\State discard the first $\mathrm{burn}=500$ days of each chain (reservoir burn-in)
\For{each chain $k=1,\dots,K$}
    \State $Y^{(k)} \gets \exp\!\big(\sum_{t=1}^{N} x_t^{(k)}\big)$ \Comment{terminal price ratio $S_{T}/S_0$, chain $k$}
\EndFor
\For{each moneyness $m$ in the ATM-weighted grid, Eq.~\eqref{eq:grid}}
    \If{$m\ge1$}
        \State $Q(m) \gets \frac1K\sum_k \max(Y^{(k)}-m,\,0)$ \Comment{OTM call, undiscounted (forward $F=1$)}
    \Else
        \State $Q(m) \gets \frac1K\sum_k \max(m-Y^{(k)},\,0)$ \Comment{OTM put}
    \EndIf
    \State $\widehat\sigma_{\text{iv}}(m) \gets$ Black-76 inversion of $Q(m)$ at $(F{=}1,K{=}m,T_{\text{years}})$ (\texttt{brentq}, bracket $[10^{-4},6.0]$)
\EndFor
\State $\sigma^2 \gets$ \texttt{cboe\_variance\_from\_grid}$(F{=}1, T_{\text{years}}, \{m\}, \{Q(m)\})$ \Comment{same routine as Section~\ref{sec:data}}
\end{algorithmic}
\end{algorithm}
Run once per expiry (near: 29d, next: 57d), $K=4{,}000$ chains, seed
base $50{,}000{,}000$, ATM-weighted moneyness grid (Eq.~\eqref{eq:grid}
above): step $0.005$ inside $[0.85,1.20]$, step $0.02$ outside, over
$[0.60,1.60]$ (104 points total). Step 12's $\sigma^2$ feeds the same 2-step CBOE
interpolation used on the market's own near/next $\sigma^2$
(Eq.~\eqref{eq:cboe}) to produce the model's own 30-day VIX
(Section~\ref{sec:results}); steps 8--11's $\widehat\sigma_{\text{iv}}(m)$
curve is the smile plotted in Section~\ref{sec:jointcalib}.

\subsection{VIX-smile algorithm: nested Monte Carlo}
\label{sec:vix-algo}
A VIX option's own payoff depends on the \emph{model-free variance
swap rate itself}, observed at a future date -- not on the underlying
index level. Reproducing its smile therefore needs a second Monte
Carlo layer nested inside the first: an \emph{outer} ensemble of
possible reservoir states on the option's own expiry date $N$, and,
\emph{from each outer draw}, an \emph{inner} ensemble of 30-day-forward
paths used to compute that day's own model-implied VIX by re-applying
Algorithm~\ref{alg:spx}'s own machinery (steps 8--12) started from that
state.

\begin{note}
\texttt{esn\_base.\_sim\_esn\_with\_params} always starts at
$r=z=0,(\sigma^{\text{used}})^2=0$ -- it cannot be conditioned on an
arbitrary mid-path state, which the inner loop requires (each of the
$K_{\text{outer}}$ outer draws has its \emph{own} state on day $N$).
A batched, vectorized reimplementation of the identical
Eqs.~\eqref{eq:reservoirupdate}--\eqref{eq:sigdef}
(\texttt{vectorized\_sim.py}), accepting an arbitrary initial state and
simulating many replicates at once, is used instead. It was validated
against the original scalar engine on 400 independent replicates,
$T=3{,}000$ (both started from $r=z=0$, matching the original's own
convention): mean, std, and $\mathrm{acf}(|x|)$ at lags 1/5/10 agree to
3 significant figures (e.g.\ $\mathrm{acf}_1=0.373$ scalar vs.\ $0.377$
vectorized); excess kurtosis differs more ($26.1$ vs.\ $20.9$ across
400 draws) but this is consistent with -- not contrary to -- the
extreme sampling variance of that particular statistic already
documented for this model.
\end{note}

\begin{algorithm}[H]
\caption{VIX-option smile at horizon $N_{\text{cal}}$ calendar days}
\label{alg:vix}
\begin{algorithmic}[1]
\State $N \gets \mathrm{round}(N_{\text{cal}}\times252/365)$
\State \textbf{Outer:} simulate $K_{\text{outer}}$ chains, $\mathrm{burn}+N$ days, from $r=z=0$; keep the terminal state $(r_N^{(k)},z_N^{(k)},(\sigma^{\text{used}}_N)^{2(k)})_{k=1}^{K_{\text{outer}}}$
\For{each outer draw $k=1,\dots,K_{\text{outer}}$}
    \State \textbf{Inner:} simulate $K_{\text{inner}}$ chains, 21 trading days ($\approx$30 calendar days), starting from state $k$
    \State $Y^{(k,j)} \gets \exp(\text{cumulative return}),\; j=1,\dots,K_{\text{inner}}$
    \State $\mathrm{VIX}_N^{(k)} \gets 100\sqrt{\sigma^2}$, with $\sigma^2\gets$ \texttt{cboe\_variance\_from\_grid} applied to $\{Y^{(k,j)}\}_j$ (Algorithm~\ref{alg:spx}, steps 8--12, ATM-weighted grid Eq.~\eqref{eq:grid}, coarser variant: step $0.01$/$0.04$, 53 points)
\EndFor
\State $F_{\text{VIX}} \gets \frac{1}{K_{\text{outer}}}\sum_k \mathrm{VIX}_N^{(k)}$ \Comment{model's own VIX future for this expiry}
\For{each real market strike $K$ in the VIX-option chain at this expiry}
    \State $Q(K) \gets \frac{1}{K_{\text{outer}}}\sum_k \max(\pm(\mathrm{VIX}_N^{(k)}-K),0)$ \Comment{OTM, sign by $K\gtrless F_{\text{VIX}}$}
    \State $\widehat\sigma_{\text{iv}}(K) \gets$ Black-76 inversion of $Q(K)$ at $(F_{\text{VIX}},K,T_{\text{years}})$
\EndFor
\end{algorithmic}
\end{algorithm}

\begin{note}
Total simulated ``chain-days'' for one expiry: outer
$K_{\text{outer}}\times(\mathrm{burn}+N)$ plus inner
$K_{\text{outer}}K_{\text{inner}}\times21$. With
$K_{\text{outer}}=K_{\text{inner}}=600$ this is
$\approx3.6\times10^5+7.6\times10^6\approx7.9\times10^6$ chain-days per
expiry, fully vectorized (batch dimension $=K_{\text{outer}}K_{\text{inner}}=360{,}000$
for the inner loop) -- run in $\approx48$s per expiry on the
container's own CPU. The 30-calendar-day window inside step 5 is
priced directly against a single 30-day horizon (no further near/next
blend) rather than the full 2-step CBOE interpolation of
Algorithm~\ref{alg:spx} -- a documented simplification, since blending
2 nested expiries per outer draw would square the cost again.
\end{note}
Seed bases: $70{,}000{,}000$ (outer, near), $+1$ (outer, next),
$+1{,}000$/$+1{,}001$ (inner). $K_{\text{outer}}=K_{\text{inner}}=600$,
$\mathrm{burn}=500$, applied at both VIX-option expiries (27d, 55d).

\subsection{Joint cross-sectional calibration (SPX + VIX)}
\label{sec:jointcalib}

\begin{equation}
\mathcal L(\theta) = \sum_{i=1}^N w_i\Big(\sigma_{\text{imp}}^\theta(K_i,T_i)-\sigma_{\text{imp}}^{\text{mkt}}(K_i,T_i)\Big)^2,
\qquad
\hat\theta = \arg\min_{\theta\in\Theta}\mathcal L(\theta).
\label{eq:calibobj}
\end{equation}
$\theta=(H,\texttt{rough\_scale},\lambda_{lo},\lambda_{hi},m_1,m_2,\alpha,\Delta b_0)$
-- every parameter the time-series joint-calibration search touches,
plus the quadratic-channel coefficient $m_2$ (Eq.~\eqref{eq:quad},
fixed at $0$ throughout the time-series calibration) and an additive
offset $\Delta b_0$ to the architecture's own bias term
(Eq.~\eqref{eq:b0}, already a native argument of
\texttt{esn\_base.\_sim\_esn\_with\_params}, no engine modification
needed) -- calibrated \emph{jointly against all 6 quoted expiries at
once}: SPX near (29d) + SPX next (57d) + VIX-option near (27d) +
VIX-option next (55d) + VIX-option 118d + VIX-option 153d.

\begin{note}
\textbf{SPX and VIX are not fit separately and then compared -- they
share 1 $\theta$.} $H$, \texttt{rough\_scale}, $\lambda_{lo}$,
$\lambda_{hi}$, $m_1$, $m_2$, $\alpha$, and $\Delta b_0$ each take a
single value that simultaneously parametrizes the SPX pricing engine
(Section~\ref{sec:spx-algo}) and the VIX pricing engine
(Section~\ref{sec:vix-algo}, itself built on top of the SPX path
simulator, since VIX is a functional of the same underlying process)
via the identical \texttt{build\_matrices\_raw(...)} call in both.
Eq.~\eqref{eq:calibobj}'s own sum runs over SPX \emph{and} VIX
residuals together in 1 objective, minimized by 1 optimizer call --
there is no separate SPX-only or VIX-only calibration step anywhere
in this note, and no mechanism by which the 2 could disagree about
$\theta$: every figure and table showing both instruments is showing
the same $\hat\theta$'s 2 different consequences, not 2 independently
tuned models placed side by side.
\end{note}

\subsubsection{Calibration target: the mid of put and call implied
vol}
\label{sec:calibtarget}
The market data carries both a put row and a call row at (almost)
every strike. Put-call parity says these should agree exactly; in
practice they agree only approximately (a few vol points apart,
growing away from the money -- visible directly in
Figures~\ref{fig:spxfinal}--\ref{fig:vixfinal}'s own put/call
markers), since each is its own noisy market print of a single
underlying implied vol. \texttt{load\_spx\_points} and
\texttt{load\_vix\_points} average the put and call mid-IV at each
strike into 1 calibration point (grouping by moneyness), rather than
treating both as independent targets: this halves the point count
(SPX near: $899\to362$) and means every point is fit on equal footing,
with no separate in-the-money/out-of-the-money convention to track.

\subsubsection{Weighting}
\label{sec:weighting}
Weights $w_i\ge0$ serve 2 purposes: correcting the heteroscedasticity
of the Black-76 price$\to$vol inversion (deep OTM options have an
ill-identified implied vol -- the price$\to$vol map is very flat
there) and weighting by vega, $w_i\propto\mathcal V(K_i,T_i)$ (Black-76
vega, from each point's own market IV, since Eq.~\eqref{eq:calibobj}
is written in \emph{vol} space). On top of vega, a Gaussian kernel
sharpens the weight toward the money, with a wider bandwidth on the
$m\ge1$ side (SPX: $0.08/0.15$; VIX: $1.5/1.5$, wide enough that VIX's
own far strikes -- routinely quoted out to $K/F\approx5$--$9$ -- still
carry real weight, not the $\approx0$ a narrower kernel gives them),
plus a flat extra multiplier on every $m\ge1$ point
(\texttt{MONEYNESS\_HIGH\_BOOST}$=3.0$), diagnosed as necessary once
direct inspection showed the model and market smiles failing to
coincide specifically on that side:
\begin{equation}
w_i \;\propto\; \mathcal V(K_i,T_i)\cdot
\exp\!\Big(-\tfrac12\big(\ln(K_i/F)/\text{bw}(K_i)\big)^2\Big)\cdot
\begin{cases}
\text{MONEYNESS\_HIGH\_BOOST} & \text{if } K_i/F\ge1 \\
1 & \text{if } K_i/F<1
\end{cases}
\label{eq:atmweight}
\end{equation}
On top of this, a cross-instrument scale normalisation divides each
expiry's own residuals by its own weighted RMSE at a fixed reference
$\theta_0$ (the $\Delta b_0$-only corrected baseline), so no expiry's
own absolute vol scale (VIX implied vols sit on a $\sim100\%+$ scale
against SPX's $\sim15$--$40\%$) dominates the objective by construction:
\begin{gather}
s_e = \sqrt{\frac{1}{\sum_{i\in e}w_i}\sum_{i\in e}w_i\big(\sigma_{\text{imp}}^{\theta_0}(K_i,T_i)-\sigma_{\text{imp}}^{\text{mkt}}(K_i,T_i)\big)^2},
\notag \\
\text{residual}_i \leftarrow \frac{\sqrt{w_i}\big(\sigma_{\text{imp}}^\theta(K_i,T_i)-\sigma_{\text{imp}}^{\text{mkt}}(K_i,T_i)\big)}{s_e},\ i\in e.
\label{eq:scalenorm}
\end{gather}
Finally, an explicit instrument-level weight
($\text{vix\_near}=3.0$, $\text{vix\_next}=1.5$) is applied to those
2 expiries' own residuals on top of Eq.~\eqref{eq:scalenorm}, since
without it the wide VIX kernel above dilutes their own near-money
weight enough that they fit distinctly worse than the other 4
expiries.

\begin{algorithm}[h]
\caption{Joint calibration}
\label{alg:jointcalib}
\begin{algorithmic}[1]
\State reparametrise $\lambda_{lo},\lambda_{hi}$ as $\log_{10}\lambda_{lo},\log_{10}\lambda_{hi}$ \Comment{both span multiple orders of magnitude; linear finite-difference steps are meaningless otherwise}
\State fix a single seed $s_0$ for every $\theta$ trial (common random numbers)
\State compute $s_e$ for each of the 6 expiries ONCE, at $\theta_0$ (Eq.~\eqref{eq:scalenorm}) \Comment{fixed for the whole run, never recomputed at other $\theta$}
\Function{Residuals}{$\theta$}
    \State build $\texttt{mat}\gets$\texttt{build\_esn\_matrices}$(H,\texttt{rough\_scale},\lambda_{lo},\lambda_{hi},m_1,m_2)$
    \For{each of the 6 expiries $e$ (SPX near/next, VIX near/next/118d/153d)}
        \If{SPX} $\sigma_{\text{imp}}^\theta(K_i)\gets$ Algorithm~\ref{alg:spx} (steps 3--14), $K_{\text{spx}}$ chains, seed $s_0$
        \Else\ (VIX) $\sigma_{\text{imp}}^\theta(K_i)\gets$ Algorithm~\ref{alg:vix} (steps 2--11), $K_{\text{outer}},K_{\text{inner}}$ chains, seed $s_0$
        \EndIf
        \State append $\text{IW}(e)\cdot\sqrt{w_i}\big(\sigma_{\text{imp}}^\theta(K_i)-\sigma_{\text{imp}}^{\text{mkt}}(K_i)\big)/s_e$ to the residual vector, all $K_i\in e$ \Comment{$\text{IW}$ = the instrument-level weight above}
    \EndFor
    \State \Return residual vector, divided by $\sqrt{6}$ so $\mathcal L(\theta_0)=1$ exactly (the usual reduced-objective convention)
\EndFunction
\State $\hat\theta\gets$ \texttt{scipy.optimize.least\_squares}(\textsc{Residuals}, $\theta_0$, bounds $\Theta$, method=\texttt{trf}, loss=\texttt{soft\_l1}), run in successive checkpointed rounds
\State \textbf{Confirm:} re-evaluate $\mathcal L(\hat\theta)$, averaged over 2 independent seeds, at $K\gg K_{\text{search}}$
\end{algorithmic}
\end{algorithm}

\begin{table}[h]
\centering
\caption{Bounds $\Theta$ and warm start $\theta_0$.}
\label{tab:calibbounds}
\begin{tabular}{lccc}
\toprule
& lower & upper & $\theta_0$ \\
\midrule
$H$ & $0.01$ & $2.5$ & $0.386$ \\
\texttt{rough\_scale} & $0.2$ & $0.7$ & $0.672$ \\
$\lambda_{lo}$ & $10^{-5}$ & $0.1$ & $10^{-5}$ \\
$\lambda_{hi}$ & $10^{-3}$ & $5.0$ & $4.881$ \\
$m_1$ & $-1.0$ & $0.5$ & $0.1745$ \\
$m_2$ & $-4.0$ & $4.0$ & $0$ \\
$\alpha$ & $0.02$ & $1.0$ & $0.220$ \\
$\Delta b_0$ & $-1.3$ & $-0.3$ & $-0.70$ \\
\bottomrule
\end{tabular}
\end{table}

\begin{note}
\textbf{Why this is computationally tractable.} A scalar,
per-path Python loop costs $\approx9$--$10$s for a single 2-expiry SPX
residual evaluation even at $K=150$, and $\approx48$s for a single
VIX-option evaluation (nested Monte Carlo) at
$K_{\text{outer}}=K_{\text{inner}}=600$ -- far too slow for a 6-expiry
objective. \texttt{vectorized\_sim.py} (Section~\ref{sec:vix-algo}),
extended with native \texttt{b0\_delta} support and used for the SPX
legs too, gives a $\approx50$--$80\times$ speed-up (validated to still
match the scalar engine, Section~\ref{sec:vix-algo}'s own Note): a
full 6-expiry residual evaluation costs $\approx15$--$20$s, not
several minutes.
\end{note}

\begin{figure}[h]
\centering
\includegraphics[width=\textwidth]{../figures/calibration_convergence.png}
\caption{$\mathcal L(\theta)/N$ (Eq.~\eqref{eq:calibobj}'s weighted
sum-of-squares, current weighting, divided by $N=858$ -- the total
number of points across all 6 specs, so the true optimum's own floor
reads as a genuinely small number rather than an arbitrary-looking
raw sum) across 10 separate optimizer runs concatenated end to end
(no marker on the plot itself distinguishes them -- explained here
instead). Left: the full range, starting from a deliberately extreme
corner of bounds $\Theta$ (Table~\ref{tab:calibbounds}) -- $H$,
$\lambda_{hi}$, $m_1$, $m_2$ all at their own upper bound, $\alpha$
and $\Delta b_0$ at their lower bound -- the single highest-loss point
anywhere in this note ($\mathcal L/N\approx1.4$, an order of magnitude
above any random draw tried elsewhere here), which converges on its
own within $\approx60$ evaluations down to $\mathcal L/N\approx0.009$.
Right: the same data, zoomed to $[0,0.06]$, where the later structure
is actually visible. Evaluations 65--160 start from a genuinely random
$\theta$ (uniform draw within $\Theta$, seed 42) and get stuck at
$\mathcal L/N\approx0.030$--$0.039$, with $\alpha$ and $m_2$ both
pinning near their own bounds -- a real, distinctly worse local
optimum, not a slow approach to the true one. Evaluations 161--190
restart from the note's own informed $\theta_0$ and converge
immediately to $\mathcal L/N\approx0.0054$, the true optimum every
other calibration in this note reaches. Evaluations 191--244 try a
2nd random $\theta$ (seed 123): this one finds its own way to the
same true optimum unaided, just less directly. Evaluations 245--335
try a 3rd random $\theta$ (seed 7) and get stuck again, at a 3rd,
different plateau ($\mathcal L/N\approx0.009$--$0.012$). Evaluations
336--353 restart again and converge to $\mathcal L/N\approx0.0054$.
Evaluations 354--459 try a 4th random $\theta$ (seed 2024): stuck for
$>60$ evaluations at $\mathcal L/N\approx0.03$--$0.038$, then breaks
free and also finds its own way to the true optimum unaided.
Evaluations 460--477 restart once more and converge to
$\mathcal L/N\approx0.0054$ again. Evaluations 478--533 try a 5th
random $\theta$ (seed 999): stuck throughout this entire budget at
$\mathcal L/N\approx0.036$--$0.052$, never finding its way to the true
optimum. Evaluations 534--590 restart a 4th time and converge to
$\mathcal L/N\approx0.0054$ once more; the remaining evaluations spent
there do not push it any lower. A handful of \texttt{trf}'s own
finite-difference Jacobian probes reach $\mathcal L/N>0.06$, off the
top of the right panel's own range (visible on the left), and are not
shown there. Across the 6 random attempts total (including the
extreme corner) here, 3 eventually found the true optimum unaided
(though far less directly than an informed start), 3 got stuck
throughout at 3 different, distinctly worse plateaus; every
informed-start attempt converged in well under 30 evaluations, every
time, to the same $\mathcal L/N\approx0.0054$ -- and no amount of
further search once there improves on it. This is a genuine,
repeatable floor for this objective under the current model and
weighting, not a slow approach to exactly $0$: the actual reason this
note warm-starts from $\theta_0$ throughout, rather than from a random
draw.}
\label{fig:convergence}
\end{figure}

Final calibrated parameters:
\begin{gather*}
\hat\theta\approx\big(H{=}0.0369,\ \texttt{rough\_scale}{=}0.200,\
\lambda_{lo}{\approx}1.1{\times}10^{-3},\ \lambda_{hi}{\approx}1.7{\times}10^{-2},\\
m_1{=}0.260,\ m_2{=}-1.477,\ \alpha{=}0.803,\ \Delta b_0{=}-0.437\big)
\end{gather*}

\begin{table}[h]
\centering
\caption{Confirm-quality weighted RMSE at $\hat\theta$, averaged over
2 independent seeds ($K_{\text{spx}}=5{,}000$,
$K_{\text{outer}}=K_{\text{inner}}=500$ each).}
\label{tab:calibresult}
\begin{tabular}{lc}
\toprule
Expiry & weighted RMSE \\
\midrule
SPX near (29d) & $0.0288$ \\
SPX next (57d) & $0.0255$ \\
VIX near (27d) & $0.2905$ \\
VIX next (55d) & $0.0792$ \\
VIX 118d & $0.0946$ \\
VIX 153d & $0.1102$ \\
\midrule
Mean (6 expiries) & $0.1048$ \\
\bottomrule
\end{tabular}
\end{table}

\begin{note}
VIX-near's own confirm-quality number carries real seed-to-seed
noise at this $K$ (the nested VIX Monte-Carlo estimator's own
variance is large enough that a single confirm draw is not, on its
own, fully trustworthy -- Figures~\ref{fig:vixfinal} and
\ref{fig:vixdist} both show this directly: the 90\% confidence band
around the model's own mean curve, and the model's own simulated
ensemble against the market's risk-neutral density, rather than a
single realization of either). Both figures are built from multiple
independent seeds for exactly this reason, not 1.
\end{note}

Figures~\ref{fig:spxfinal} and \ref{fig:vixfinal} show the fitted
smile against the market at every quoted maturity.

\begin{figure}[h]
\centering
\includegraphics[width=0.85\textwidth]{../figures/spx_smile_final.png}
\caption{SPX smile, market vs.\ $\hat\theta$, 4 maturities: the 2
calibrated against (29d, 57d) plus 2 held out entirely from
calibration (119d, 211d). Open grey circles and grey crosses are the
market's own mid implied vol (average of bid/ask) at each traded
strike, from the put and call quote respectively -- both are shown
for transparency, though $\hat\theta$ itself is fit against their own
average, 1 point per strike (Section~\ref{sec:calibtarget}), since
they are 2 noisy prints of 1 underlying implied vol, not 2
independent facts. This is not a bid/ask spread (each marker is
already its own bid/ask mid) -- their own small disagreement (a few
vol points, growing away from the money) is a genuine feature of the
quoted data. The blue curve is the \emph{mean} model IV across 10
independent Monte-Carlo seeds, with a shaded 90\% interval (5th--95th
percentile) -- not a single realization, since VIX pricing noise
(Figure~\ref{fig:vixfinal}) makes single draws unreliable and SPX is
shown the same way for consistency. Both the axis range and the
curve are limited to each panel's own market-quoted moneyness range.
The band stays narrow throughout: SPX pricing is cheap enough
(Section~\ref{sec:jointcalib}'s own vectorized engine) that
$K=3{,}000$ per seed leaves little residual MC noise.}
\label{fig:spxfinal}
\end{figure}

\begin{figure}[h]
\centering
\includegraphics[width=0.85\textwidth]{../figures/vix_smile_final.png}
\caption{VIX-option smile, market vs.\ $\hat\theta$, all 4 maturities
the VIX-option chain quotes, all genuinely calibrated against. As in
Figure~\ref{fig:spxfinal}, open circles and crosses are the market's
own put and call mid implied vol, and the blue curve is the
\emph{mean} model IV across 5 independent Monte-Carlo seeds with a
shaded 90\% interval. The band is visibly wide here -- unlike SPX,
VIX pricing needs the much more expensive nested Monte Carlo
(Section~\ref{sec:vix-algo}), so 5 seeds at a search-affordable $K$
leave real, visible residual noise. The jumps visible in the mean
curve at the largest strikes (most obvious in the 55d and 153d
panels) are not a plotting bug: at those strikes some, but not all,
of the 5 seeds hit the finite-sample precision limit described below
and return no value, so the mean is computed over a shrinking,
discontinuously changing set of seeds as moneyness increases -- a
visual indicator of exactly where the curve stops being trustworthy.
Both the axis range and the curve are limited to each panel's own
market-quoted moneyness range -- but unlike the SPX panels, the model
curve here often stops \emph{before} the axis' own right edge: VIX
strikes are quoted out to moneyness $\approx5$--$9$, far into the
tail of the model's own Monte-Carlo ensemble (Section~\ref{sec:results}),
where the price is indistinguishable from $0$ and the Black-76
inversion has nothing left to invert -- a genuine finite-sample
precision limit (at $K/F=9.2$, pricing a call requires simulated
$\mathrm{VIX}_N$ draws exceeding $9.2\times F_{\mathrm{VIX}}\approx160$,
and essentially none of the model's own 1,000 draws reach that level),
not an arbitrary cutoff or a fitting failure: it is a direct,
quantitative sign that the model's own risk-neutral density assigns
near-zero probability to a VIX spike that large. The same mechanism
operates during calibration itself: Eq.~\eqref{eq:calibobj}'s residual
is set to $0$ wherever the model returns no value, so any market quote
beyond the model's reach contributes no gradient to $\hat\theta$ --
silently excluded, not down-weighted. 27d and 153d both still show a
real, visible gap even under this fully corrected objective -- 27d's
model curve sits above the market throughout, and 153d's mean curve
tracks reasonably to $K/F\approx3$ before falling away -- a residual
limitation of the shared 8-parameter architecture fitting 4 VIX
maturities' own term structure at once, not of the objective.}
\label{fig:vixfinal}
\end{figure}

\subsubsection{Generalization to held-out SPX maturities}
\label{sec:heldout}
Figure~\ref{fig:spxfinal} also shows 2 SPX maturities (119d, 211d)
included in neither Eq.~\eqref{eq:calibobj}'s sum nor the scale
normalisation of Eq.~\eqref{eq:scalenorm}, priced at the same final
$\hat\theta$ purely to see whether the fit generalizes. It does not,
uniformly: at both maturities the model's put wing ($K/F<1$) sits
\emph{below} the market (undershooting, whereas the 2 calibrated
maturities track the market closely there), growing with maturity --
consistent with a genuine term-structure mismatch the 2-maturity SPX
calibration set cannot see, rather than noise. This is expected, not
a contradiction of anything claimed above: $\hat\theta$ was never
asked to fit these 2 points, and a $\theta$ fit to 2 near-dated SPX
expiries has no particular reason to extrapolate correctly to
4--7$\times$ longer maturities. VIX has no equivalent held-out
maturity in this note: the same generalization gap, diagnosed at 118d
and 153d, was addressed directly by adding both to the calibration
set (Section~\ref{sec:calibtarget}), which is why all 4 VIX panels in
Figure~\ref{fig:vixfinal} are marked ``calibrated''.

\FloatBarrier
\section{Test 1: VIX term structure}
\label{sec:results}

\begin{table}[h]
\centering
\caption{Market vs.\ model, 2026-02-19.}
\label{tab:vixresult}
\begin{tabular}{lcc}
\toprule
& Market & ESN (SP500) \\
\midrule
Near-term implied vol (29d) & 19.55\% & 35.99\% \\
Next-term implied vol (57d) & 20.76\% & 36.10\% \\
30-day VIX (interpolated) & 19.63 & 36.00 \\
\bottomrule
\end{tabular}
\end{table}

\begin{figure}[h]
\centering
\includegraphics[width=\textwidth]{../figures/vix_distribution_comparison.png}
\caption{VIX empirical distribution, all 4 quoted maturities, all
genuinely calibrated against (Section~\ref{sec:jointcalib}): the
market's own
risk-neutral density
(grey curve, Breeden-Litzenberger -- $f_Y(K)=\partial^2C(K)/\partial K^2$)
vs.\ the model's own $\hat\theta$ ensemble (blue histogram, 700
simulated draws of $\mathrm{VIX}_N$ from Algorithm~\ref{alg:vix}'s own
outer loop). Unlike Table~\ref{tab:vixresult} above (the original,
uncalibrated $\theta$), this uses the fully calibrated, final
$\hat\theta$. The market density is built from a
degree-5 polynomial fit to the chain's own mid-IV quotes in
log-moneyness, priced from a single smooth branch (calls throughout
the whole strike range, not the usual put-below/call-above OTM
splice) and differentiated analytically via a cubic-spline
interpolant -- both steps needed: the OTM splice's own first
derivative has a genuine jump of exactly 1 at the crossover strike
(differentiating the put-call parity identity $C-P=F-K$), which is a
real kink, not market noise, and produces a spurious spike under naive
finite differencing.}
\label{fig:vixdist}
\end{figure}

The model overstates every point on the term structure by a near-constant
factor of $\approx1.8\times$. This is the direct cross-sectional
counterpart of the Note in Section~\ref{sec:feedback}: $m_1>0$ biases
the unconditional variance level upward, and since $\theta$ was never
fit against any option-implied target, there is no reason to expect
this factor to be close to 1. Once $\hat\theta$ \emph{is} fit
(Section~\ref{sec:jointcalib}), Figure~\ref{fig:vixdist} shows the
level gap has largely closed on its own -- the model's own ensemble
mean sits at 15.8 (27d), 15.4 (55d), 15.9 (118d), and 15.6 (153d),
close to the market's own forward at every maturity (20.1, 20.7, 21.2,
21.7) -- even though absolute level was never a direct calibration
target, only the smile's own shape was. What remains is a genuine
shape
difference, present at all 4 maturities: the model's distribution is
visibly narrower and more concentrated than the market's, whose own
Breeden-Litzenberger density has a longer, heavier right tail out
past $\mathrm{VIX}=40$. This is the same
finding as Figures~\ref{fig:spxfinal}--\ref{fig:vixfinal}'s own smile
comparison, now shown directly on the underlying distribution rather
than through option prices: the model's continuous, jump-free dynamics
structurally under-produce the probability of a genuine volatility
spike that the market's own risk-neutral density still prices in.

\FloatBarrier
\section{Summary}
\label{sec:summary}
As calibrated purely on time-series shape statistics, the SP500
feedback-layer model gets several genuine, non-trivial cross-sectional
features right (a real SPX put skew; a real, upward-sloping VIX skew)
while getting \emph{level} wrong throughout ($\approx1.8\times$ too
high on both surfaces) and showing a real shape mismatch on both
smiles. A genuine joint recalibration against all 6 quoted expiries
at once (Section~\ref{sec:jointcalib}) -- SPX near/next, VIX
near/next/118d/153d, calibrated on the mid of each strike's put and
call implied vol, with vega-weighted, ATM-sharpened, moneyness-boosted
weights and a cross-instrument scale normalisation -- closes most of
this gap. The final $\hat\theta$ gives confirm-quality weighted RMSE
of $0.029$/$0.026$ on SPX near/next and $0.29$/$0.08$/$0.09$/$0.11$ on
VIX near/next/118d/153d (Table~\ref{tab:calibresult}), with the SPX
smile's own $m\ge1$/$m<1$ RMSE ratio -- the direct measure of whether
the model and market coincide away from the money -- improved from
$\approx3.5\times$ to $\approx1.4\times$. The VIX smile tracks the
market closely across most of its quoted range at 55d and 118d, with
a smaller residual gap at 27d and 153d: a genuine limitation of the
shared 8-parameter architecture fitting 4 VIX maturities' own term
structure at once, not a fitting failure at any 1 of them. On the SPX
side, 2 entirely held-out maturities (119d, 211d, Section~\ref{sec:heldout})
show the fit does not generalize uniformly beyond the 2 calibrated
expiries: the put wing undershoots increasingly with maturity, an
honest limit of a 2-maturity calibration set. None of the original
gaps contradicts the note family's own reported time-series results
-- none of level, SPX shape, or VIX shape was ever a time-series
calibration target -- and the recalibration shows those gaps are
addressable, with the specific methodology now documented and
reusable.

\section*{Reproducibility}
Delivered alongside this note, all built directly on the uploaded,
unmodified engine files:
\begin{itemize}
\item \texttt{market\_vix.py} -- CBOE VIX/VVIX replication, exposing
the shared variance-from-grid core function
\item \texttt{model\_pricing.py} -- Algorithm~\ref{alg:spx}
\item \texttt{vectorized\_sim.py} and its validation script -- the
batched, state-conditionable engine used throughout
Section~\ref{sec:pricing}, now with native \texttt{b0\_delta} support,
validated against the original scalar engine
\item \texttt{vix\_smile.py} -- Algorithm~\ref{alg:vix}
\item \texttt{calibrate\_smile.py} -- Section~\ref{sec:jointcalib},
Algorithm~\ref{alg:jointcalib}; includes checkpointing and a
per-evaluation objective-history log, so a long optimizer run can be
inspected, plotted, or resumed
\item \texttt{vix\_distribution.py} -- Section~\ref{sec:results}'s
Breeden-Litzenberger market density and model ensemble,
Figure~\ref{fig:vixdist}
\item \texttt{make\_figures.py} -- regenerates every figure in this
note from the JSON results the other scripts produce
\end{itemize}
All seeds are fixed and stated in
Sections~\ref{sec:spx-algo}--\ref{sec:jointcalib}.

\end{document}
